\documentclass[10pt]{beamer}
\usepackage[utf8]{inputenc}
\usepackage[T1]{fontenc}
\usepackage{lmodern}
\usepackage{amsmath}
\usepackage{amssymb}
\usepackage{listings}
\usepackage{xcolor}
\usepackage{booktabs}
\usepackage{array}
\usepackage{tabularx}
\usepackage{graphicx}
\usepackage{eurosym}
\usepackage{tikz}

\usetheme{Madrid}
\usecolortheme{default}

% SQL / code listing style
\lstdefinestyle{sqlstyle}{
  language=SQL,
  basicstyle=\small\ttfamily,
  keywordstyle=\color{blue!80!black}\bfseries,
  commentstyle=\color{gray!70}\itshape,
  stringstyle=\color{orange!80!black},
  breaklines=true,
  showstringspaces=false,
  frame=single,
  backgroundcolor=\color{gray!10},
  rulecolor=\color{gray!40},
  tabsize=2,
}

% Relational algebra shorthand commands
\newcommand{\rasel}{\sigma}
\newcommand{\raproj}{\pi}
\newcommand{\raren}{\rho}
\newcommand{\rajoin}{\bowtie}
\newcommand{\ralsem}{\ltimes}
\newcommand{\rarsem}{\rtimes}

\title[Databases I -- Exercises]{Databases I -- Exercises}
\subtitle{Semester 3 (Theory)}
\author[Neunzig, M.]{Mathis Neunzig}
\institute[DHBW Mosbach]{DHBW Mosbach}
\date[Autumn 2026]{Autumn 2026}

\AtBeginSection[]{
  \begin{frame}
    \frametitle{Table of Contents}
    \tableofcontents[currentsection]
  \end{frame}
}

\begin{document}

\frame{\titlepage}

\begin{frame}
  \frametitle{Table of Contents}
  \tableofcontents
\end{frame}

% ============================================================
\section{Relational Databases -- Keys}
% ============================================================

% ------------------------------------------------------------
\subsection{Exercise 1A -- University Schema}

\begin{frame}{Exercise 1A -- University Schema}
  \begin{block}{Relational Schema}
    \small
    \begin{itemize}
      \item \textbf{Student}(\underline{MatrNr}, FirstName, LastName, Email, DOB, MajorID\textsuperscript{FK})
      \item \textbf{Department}(\underline{DeptID}, Name, Building, HeadProfID\textsuperscript{FK})
      \item \textbf{Course}(\underline{CourseID}, Title, Credits, DeptID\textsuperscript{FK})
      \item \textbf{Enrollment}(\underline{MatrNr\textsuperscript{FK}, CourseID\textsuperscript{FK}}, Semester, Grade)
      \item \textbf{Professor}(\underline{ProfID}, FirstName, LastName, Email, DeptID\textsuperscript{FK})
    \end{itemize}
  \end{block}

  \begin{block}{Tasks -- answer for every relation}
    \begin{enumerate}
      \item List all \textbf{attributes} and their domains.
      \item Identify all \textbf{candidate keys}; select and justify your \textbf{primary key}.
      \item Identify all \textbf{foreign keys} and state which PK each one references.
      \item Classify every key as \emph{simple/composite} and \emph{natural/surrogate}.
      \item Create \textbf{2 realistic sample tuples} for each relation.
    \end{enumerate}
  \end{block}

  \begin{alertblock}{Difficulty: Simple \quad Time: 15 min}
  \end{alertblock}
\end{frame}

% ------------------------------------------------------------
\subsection{Exercise 1B -- Airline Schema}

\begin{frame}{Exercise 1B -- Airline Schema}
  \begin{block}{Relational Schema}
    \small
    \begin{itemize}
      \item \textbf{Airline}(\underline{AirlineCode}, Name, Country, FoundedYear)
      \item \textbf{Airport}(\underline{IATA}, Name, City, Country, Terminals)
      \item \textbf{Connection}(\underline{ConnectionID}, DepartureIATA\textsuperscript{FK}, ArrivalIATA\textsuperscript{FK}, AirlineCode\textsuperscript{FK}, FlightNumber)
      \item \textbf{Schedule}(\underline{ConnectionID\textsuperscript{FK}, Date}, DepartureTime, ArrivalTime, Status)
      \item \textbf{Passenger}(\underline{PassengerID}, FirstName, LastName, Passport, Nationality)
    \end{itemize}
  \end{block}

  \begin{block}{Tasks -- same 5 questions as 1A}
    \begin{enumerate}
      \item List all \textbf{attributes} and their domains.
      \item Identify all \textbf{candidate keys}; select and justify your \textbf{primary key}.
      \item Identify all \textbf{foreign keys} and state which PK each one references.
      \item Classify every key as \emph{simple/composite} and \emph{natural/surrogate}.
      \item Create \textbf{2 realistic sample tuples} for each relation.
    \end{enumerate}
  \end{block}

  \begin{alertblock}{Difficulty: Simple \quad Time: 15 min}
  \end{alertblock}
\end{frame}

% ------------------------------------------------------------
\subsection{Exercise 1C -- Event Schema}

\begin{frame}{Exercise 1C -- Event Schema}
  \begin{block}{Relational Schema}
    \small
    \begin{itemize}
      \item \textbf{Event}(\underline{EventID}, Title, Date, Location, OrganizerID\textsuperscript{FK})
      \item \textbf{Organizer}(\underline{OrgID}, Name, Email, Phone)
      \item \textbf{Sponsor}(\underline{SponsorID}, CompanyName, ContactEmail, Budget)
      \item \textbf{Sponsorship}(\underline{EventID\textsuperscript{FK}, SponsorID\textsuperscript{FK}}, Amount, ContractDate)
      \item \textbf{Attendee}(\underline{AttendeeID}, FirstName, LastName, Email)
      \item \textbf{Attendance}(\underline{AttendeeID\textsuperscript{FK}, EventID\textsuperscript{FK}}, RegistrationDate, Paid)
    \end{itemize}
  \end{block}

  \begin{block}{Tasks -- same 5 questions as 1A}
    \begin{enumerate}
      \item List all \textbf{attributes} and their domains.
      \item Identify all \textbf{candidate keys}; select and justify your \textbf{primary key}.
      \item Identify all \textbf{foreign keys} and state which PK each one references.
      \item Classify every key as \emph{simple/composite} and \emph{natural/surrogate}.
      \item Create \textbf{2 realistic sample tuples} for each relation.
    \end{enumerate}
  \end{block}

  \begin{alertblock}{Difficulty: Intermediate \quad Time: 20 min}
  \end{alertblock}
\end{frame}

% ------------------------------------------------------------
\subsection{Exercise 1D -- Music Schema}

\begin{frame}{Exercise 1D -- Music Schema}
  \begin{block}{Relational Schema}
    \small
    \begin{itemize}
      \item \textbf{Artist}(\underline{ArtistID}, Name, Country, DebutYear, Genre)
      \item \textbf{Album}(\underline{AlbumID}, Title, ArtistID\textsuperscript{FK}, ReleaseYear, Label)
      \item \textbf{Track}(\underline{TrackID}, Title, AlbumID\textsuperscript{FK}, TrackNumber, Duration)
      \item \textbf{Features}(\underline{TrackID\textsuperscript{FK}, FeaturedArtistID\textsuperscript{FK}})
      \item \textbf{Playlist}(\underline{PlaylistID}, Name, UserID, CreatedDate)
      \item \textbf{PlaylistTrack}(\underline{PlaylistID\textsuperscript{FK}, TrackID\textsuperscript{FK}}, AddedDate, Position)
    \end{itemize}
  \end{block}

  \begin{block}{Tasks -- same 5 questions as 1A}
    \begin{enumerate}
      \item List all \textbf{attributes} and their domains.
      \item Identify all \textbf{candidate keys}; select and justify your \textbf{primary key}.
      \item Identify all \textbf{foreign keys} and state which PK each one references.
      \item Classify every key as \emph{simple/composite} and \emph{natural/surrogate}.
      \item Create \textbf{2 realistic sample tuples} for each relation.
    \end{enumerate}
  \end{block}

  \begin{alertblock}{Difficulty: Intermediate \quad Time: 20 min}
  \end{alertblock}
\end{frame}

% ============================================================
\section{Relational Algebra}
% ============================================================

% ------------------------------------------------------------
\subsection{Exercise 2A -- Customer / Supplier / Good (Basic)}

\begin{frame}{Exercise 2A -- Schema}
  \begin{block}{Schema}
    \begin{itemize}
      \item \textbf{Customer}(\underline{CustID}, Name, City, Balance)
      \item \textbf{Supplier}(\underline{SupplierID}, Name, City, ContactEmail)
      \item \textbf{Good}(\underline{GoodID}, Name, Price, SupplierID\textsuperscript{FK})
    \end{itemize}
  \end{block}

  \vspace{4pt}
  \small
  \begin{tabular}{|l|l|l|r|}
    \hline
    \textbf{CustID} & \textbf{Name} & \textbf{City} & \textbf{Balance} \\
    \hline
    C1 & Maier GmbH & Mosbach & 240.00 \\
    C2 & Schulz AG & Heidelberg & -50.00 \\
    \hline
  \end{tabular}
  \quad
  \begin{tabular}{|l|l|l|}
    \hline
    \textbf{SupplierID} & \textbf{Name} & \textbf{City} \\
    \hline
    S1 & Fresh Foods & Mosbach \\
    S2 & Top Drinks & Heidelberg \\
    \hline
  \end{tabular}

  \vspace{4pt}
  \begin{tabular}{|l|l|r|l|}
    \hline
    \textbf{GoodID} & \textbf{Name} & \textbf{Price} & \textbf{SupplierID} \\
    \hline
    G1 & Apple & 0.49 & S1 \\
    G2 & Orange Juice & 2.99 & S2 \\
    G3 & Sparkling Water & 6.50 & S1 \\
    \hline
  \end{tabular}
\end{frame}

\begin{frame}{Exercise 2A -- Tasks}
  \begin{block}{Tasks -- express in relational algebra}
    \begin{enumerate}
      \item List all goods with price $> 5.00$ (show Name, Price).
      \item List all suppliers located in Mosbach.
      \item List all customers with a \textbf{negative} balance.
      \item Find all goods supplied by suppliers in Mosbach.
      \item List customer names and cities only (project exactly these two columns).
      \item Find customers whose balance is less than the price of the cheapest good.
    \end{enumerate}
  \end{block}

  \begin{block}{Reminder: operators}
    $\sigma_{\text{condition}}(R)$ \quad
    $\pi_{\text{attrs}}(R)$ \quad
    $R_1 \bowtie R_2$ \quad
    $R_1 \times R_2$ \quad
    $R_1 \setminus R_2$
  \end{block}

  \begin{alertblock}{Difficulty: Simple \quad Time: 15 min}
  \end{alertblock}
\end{frame}

% ------------------------------------------------------------
\subsection{Exercise 2B -- Customer / Supplier / Good (Harder)}

\begin{frame}{Exercise 2B -- Tasks}
  \begin{block}{Schema (same as 2A, Supplier now also has a Cost column)}
    \begin{itemize}
      \item \textbf{Customer}(\underline{CustID}, Name, City, Balance)
      \item \textbf{Supplier}(\underline{SupplierID}, Name, City, ContactEmail, Cost)
      \item \textbf{Good}(\underline{GoodID}, Name, Price, SupplierID\textsuperscript{FK})
    \end{itemize}
  \end{block}

  \begin{block}{Tasks}
    \begin{enumerate}
      \item Find goods where the price margin (Price $-$ Cost) $> 2.00$.
      \item List all suppliers in Heidelberg.
      \item List customers whose balance exceeds the price of the most expensive good.
      \item Find all goods \textbf{not} supplied by suppliers in Mosbach.
      \item Find all customer--supplier pairs from the \textbf{same city} (use Cartesian product + restriction).
    \end{enumerate}
  \end{block}

  \begin{alertblock}{Difficulty: Simple--Intermediate \quad Time: 20 min}
  \end{alertblock}
\end{frame}

% ------------------------------------------------------------
\subsection{Exercise 2C -- PC / Laptop / Printer (Intermediate)}

\begin{frame}{Exercise 2C -- Schema}
  \begin{block}{Schema}
    \begin{itemize}
      \item \textbf{Product}(\underline{ProductID}, Maker, Type)
      \item \textbf{PC}(\underline{ProductID\textsuperscript{FK}}, Speed\_GHz, RAM\_GB, HDD\_GB, Price)
      \item \textbf{Laptop}(\underline{ProductID\textsuperscript{FK}}, Speed\_GHz, RAM\_GB, HDD\_GB, Screen\_inch, Price)
      \item \textbf{Printer}(\underline{ProductID\textsuperscript{FK}}, Color, Type, Price)
    \end{itemize}
  \end{block}

  \small
  \begin{tabular}{|l|l|l|}
    \hline
    \textbf{ProductID} & \textbf{Maker} & \textbf{Type} \\
    \hline
    P01 & Dell & PC \\
    P02 & Dell & Laptop \\
    P03 & HP & PC \\
    P04 & HP & Printer \\
    P05 & Apple & Laptop \\
    \hline
  \end{tabular}
  \quad
  \begin{tabular}{|l|r|r|r|r|}
    \hline
    \textbf{ProductID} & \textbf{GHz} & \textbf{RAM} & \textbf{HDD} & \textbf{Price} \\
    \hline
    P01 & 3.2 & 16 & 512 & 899.00 \\
    P03 & 2.8 & 8 & 256 & 649.00 \\
    \hline
  \end{tabular}
\end{frame}

\begin{frame}{Exercise 2C -- Tasks}
  \begin{block}{Tasks}
    \begin{enumerate}
      \item List all PCs with speed $> 3.0$ GHz.
      \item Find all manufacturers who make \textbf{both} PCs and Laptops.
      \item List all products from manufacturer \texttt{'Dell'}.
      \item Find all PC--Laptop pairs from the \textbf{same manufacturer} (use Cartesian product).
      \item Find makers who produce at least \textbf{3 different types} of products.
    \end{enumerate}
  \end{block}

  \begin{block}{Hint for task 2}
    Use projection + intersection:
    $\pi_{\text{Maker}}(\sigma_{\text{Type}=\text{'PC'}}(\text{Product})) \cap \pi_{\text{Maker}}(\sigma_{\text{Type}=\text{'Laptop'}}(\text{Product}))$
  \end{block}

  \begin{alertblock}{Difficulty: Intermediate \quad Time: 20 min}
  \end{alertblock}
\end{frame}

% ------------------------------------------------------------
\subsection{Exercise 2D -- PC / Laptop / Printer (Advanced Algebra)}

\begin{frame}{Exercise 2D -- Tasks}
  \begin{block}{Schema (same as 2C)}
    \textbf{Product}(\underline{ProductID}, Maker, Type) \quad
    \textbf{PC}(\underline{ProductID}, Speed\_GHz, RAM\_GB, HDD\_GB, Price)\\
    \textbf{Laptop}(\underline{ProductID}, Speed\_GHz, RAM\_GB, HDD\_GB, Screen\_inch, Price)\\
    \textbf{Printer}(\underline{ProductID}, Color, Type, Price)
  \end{block}

  \begin{block}{Tasks}
    \begin{enumerate}
      \item Find laptops with screen size $\geq 15$ inch \textbf{and} price $< 1000$.
      \item List all manufacturers who \textbf{only} make printers (not PCs or Laptops).
      \item Find all products more expensive than the \textbf{average} PC price (use $\gamma$).
      \item Count products per manufacturer using $\gamma$ (grouping/aggregation).
      \item Write as relational algebra: \emph{``Find makers of PCs with speed $> 3$ GHz''} -- express as a single nested expression.
    \end{enumerate}
  \end{block}

  \begin{alertblock}{Difficulty: Intermediate \quad Time: 25 min}
  \end{alertblock}
\end{frame}

% ============================================================
\section{Normal Forms (1NF -- BCNF)}
% ============================================================

% ------------------------------------------------------------
\subsection{Exercise 3A -- Eurovision Song Contest (Worked Example)}

\begin{frame}{Exercise 3A -- Eurovision Song Contest}
  \begin{block}{Starting Relation}
    $R($Year, Country, Song, Artist, Points, Rank, Host\_Country, Host\_City$)$
  \end{block}

  \small
  \begin{tabular}{|l|l|l|l|r|r|l|l|}
    \hline
    \textbf{Year} & \textbf{Country} & \textbf{Song} & \textbf{Artist} & \textbf{Pts} & \textbf{Rank} & \textbf{Host} & \textbf{City} \\
    \hline
    2023 & Sweden  & Tattoo      & Loreen           & 583 & 1 & UK    & Liverpool \\
    2023 & Finland & Cha Cha Cha & K\"{a}\"{a}rij\"{a}           & 526 & 2 & UK    & Liverpool \\
    2022 & Ukraine & Stefania    & Kalush Orchestra & 631 & 1 & Italy & Turin     \\
    2022 & UK      & Space Man   & Sam Ryder        & 466 & 2 & Italy & Turin     \\
    \hline
  \end{tabular}

  \vspace{6pt}
  \begin{block}{Tasks}
    \begin{enumerate}
      \item Identify all violations of 1NF, 2NF, 3NF, BCNF.
      \item Decompose to \textbf{1NF} (if needed).
      \item Decompose to \textbf{2NF} -- show functional dependencies, show decomposed tables.
      \item Decompose to \textbf{3NF} -- identify and eliminate transitive FDs.
      \item Check if \textbf{BCNF} is achieved; decompose further if necessary.
    \end{enumerate}
  \end{block}
\end{frame}

\begin{frame}{Exercise 3A -- Functional Dependencies}
  \begin{block}{Identify the Functional Dependencies}
    Working from the relation\\
    $R($Year, Country, Song, Artist, Points, Rank, Host\_Country, Host\_City$)$:
    \vspace{4pt}
    \begin{itemize}
      \item Which attributes uniquely identify a tuple?
      \item What does \{Year, Country\} determine?
      \item What does Year alone determine?
      \item Is there a transitive dependency (e.g.\ $A \to B \to C$)?
    \end{itemize}
  \end{block}

  \begin{block}{Guiding Questions for Each Normal Form}
    \begin{description}
      \item[1NF] Are all attribute values atomic?  Any repeating groups?
      \item[2NF] Are there non-prime attributes that depend on only \emph{part} of the PK?
      \item[3NF] Are there non-prime attributes that transitively depend on the PK?
      \item[BCNF] Does every non-trivial FD have a \textbf{superkey} on the left?
    \end{description}
  \end{block}

  \begin{alertblock}{Difficulty: Simple--Intermediate \quad Time: 25 min}
  \end{alertblock}
\end{frame}

% ------------------------------------------------------------
\subsection{Exercise 3B -- Ticket Schema}

\begin{frame}{Exercise 3B -- Ticket Schema}
  \begin{block}{Starting Relation}
    $\text{Ticket}($TicketID, EventName, EventDate, Venue, VenueCity, CustomerName,\\
    \hspace{2.4cm} CustomerEmail, SeatNumber, Price, TicketType$)$
  \end{block}

  \begin{block}{Known Functional Dependencies}
    \begin{itemize}
      \item $\text{TicketID} \to \text{all other attributes}$
      \item $\text{EventName} \to \text{EventDate},\ \text{Venue}$
      \item $\text{Venue} \to \text{VenueCity}$
      \item $\text{CustomerEmail} \to \text{CustomerName}$
      \item $\{\text{TicketID},\ \text{SeatNumber}\} \to \text{Price}$ \quad (if modelled as composite)
    \end{itemize}
  \end{block}

  \begin{block}{Tasks}
    \begin{enumerate}
      \item Identify all violations of 1NF, 2NF, 3NF, BCNF.
      \item Decompose step by step to 1NF $\to$ 2NF $\to$ 3NF.
      \item Show all intermediate tables after each decomposition step.
      \item Check if BCNF is achieved.
    \end{enumerate}
  \end{block}

  \begin{alertblock}{Difficulty: Intermediate \quad Time: 30 min}
  \end{alertblock}
\end{frame}

% ------------------------------------------------------------
\subsection{Exercise 3C -- Library Loan Schema}

\begin{frame}{Exercise 3C -- Library Loan Schema}
  \begin{block}{Starting Relation}
    $\text{LibraryLoan}($LoanID, BookISBN, BookTitle, AuthorName, BorrowerID,\\
    \hspace{3.0cm} BorrowerName, BorrowerEmail, LoanDate, ReturnDate,\\
    \hspace{3.0cm} LibraryBranch, BranchCity$)$
  \end{block}

  \begin{block}{Known Functional Dependencies}
    \begin{itemize}
      \item $\text{LoanID} \to \text{all other attributes}$
      \item $\text{BookISBN} \to \text{BookTitle},\ \text{AuthorName}$
      \item $\text{BorrowerID} \to \text{BorrowerName},\ \text{BorrowerEmail}$
      \item $\text{LibraryBranch} \to \text{BranchCity}$
    \end{itemize}
  \end{block}

  \begin{block}{Tasks}
    \begin{enumerate}
      \item Identify all violations of 1NF, 2NF, 3NF, BCNF.
      \item Decompose step by step to 1NF $\to$ 2NF $\to$ 3NF.
      \item Show all intermediate tables after each decomposition step.
      \item Verify lossless-join property after final decomposition.
    \end{enumerate}
  \end{block}

  \begin{alertblock}{Difficulty: Advanced \quad Time: 35 min}
  \end{alertblock}
\end{frame}

% ============================================================
\section{Normal Forms (4NF / 5NF)}
% ============================================================

% ------------------------------------------------------------
\subsection{Exercise 4A -- Train / Conductor / Mechanic (4NF)}

\begin{frame}{Exercise 4A -- Train / Conductor / Mechanic}
  \begin{block}{Starting Relation}
    $R($Train, Conductor, Mechanic$)$
  \end{block}

  \begin{block}{Multi-Valued Dependencies (MVDs)}
    \begin{itemize}
      \item $\text{Train} \twoheadrightarrow \text{Conductor}$ \quad (a train can be driven by multiple conductors, independently of mechanics)
      \item $\text{Train} \twoheadrightarrow \text{Mechanic}$ \quad (a train can be maintained by multiple mechanics, independently)
    \end{itemize}
  \end{block}

  \small
  \begin{tabular}{|l|l|l|}
    \hline
    \textbf{Train} & \textbf{Conductor} & \textbf{Mechanic} \\
    \hline
    ICE 42 & Hoffmann & Bauer \\
    ICE 42 & Hoffmann & Klein \\
    ICE 42 & Neumann  & Bauer \\
    ICE 42 & Neumann  & Klein \\
    \hline
  \end{tabular}

  \vspace{6pt}
  \begin{block}{Tasks}
    \begin{enumerate}
      \item Identify the MVDs and explain the resulting redundancy.
      \item Decompose to \textbf{4NF} -- show the decomposed tables.
      \item Verify that the decomposition is lossless.
    \end{enumerate}
  \end{block}

  \begin{alertblock}{Difficulty: Intermediate \quad Time: 20 min}
  \end{alertblock}
\end{frame}

% ------------------------------------------------------------
\subsection{Exercise 4B -- Line / Station / Day (5NF)}

\begin{frame}{Exercise 4B -- Line / Station / Day}
  \begin{block}{Starting Relation}
    $R($Line, Station, Day$)$
  \end{block}

  \begin{block}{Business Constraint (Join Dependency)}
    If line $L$ stops at station $S$, and line $L$ runs on day $D$, and station $S$ is served on day $D$ -- then the combination $(L, S, D)$ \textbf{must} exist in $R$.
  \end{block}

  \small
  \begin{tabular}{|l|l|l|}
    \hline
    \textbf{Line} & \textbf{Station} & \textbf{Day} \\
    \hline
    RE7 & Mosbach & Monday \\
    RE7 & Heidelberg & Monday \\
    RE7 & Mosbach & Friday \\
    RE7 & Heidelberg & Friday \\
    \hline
  \end{tabular}

  \vspace{6pt}
  \begin{block}{Tasks}
    \begin{enumerate}
      \item Identify the \textbf{join dependency} and explain why 4NF is not sufficient.
      \item Decompose to \textbf{5NF} -- show the three binary projections.
      \item Reconstruct the original relation by joining the projections and verify no information is lost.
    \end{enumerate}
  \end{block}

  \begin{alertblock}{Difficulty: Intermediate \quad Time: 20 min}
  \end{alertblock}
\end{frame}

% ------------------------------------------------------------
\subsection{Exercise 4C -- Train / Route / Service / CarriageType (4NF)}

\begin{frame}{Exercise 4C -- Train / Route / Service / CarriageType}
  \begin{block}{Starting Relation}
    $R($Train, Route, Service, CarriageType$)$
  \end{block}

  \begin{block}{Multi-Valued Dependencies}
    All three MVDs hold independently:
    \begin{itemize}
      \item $\text{Train} \twoheadrightarrow \text{Route}$
      \item $\text{Train} \twoheadrightarrow \text{Service}$
      \item $\text{Train} \twoheadrightarrow \text{CarriageType}$
    \end{itemize}
    (A train serves multiple routes, offers multiple services, and can have multiple carriage types -- all choices are independent of each other.)
  \end{block}

  \begin{block}{Tasks}
    \begin{enumerate}
      \item Explain why storing all combinations leads to redundancy.
      \item Decompose to \textbf{4NF} by eliminating all non-trivial MVDs.
      \item How many tables result?  Write out each table schema.
      \item Verify lossless-join property.
    \end{enumerate}
  \end{block}

  \begin{alertblock}{Difficulty: Intermediate \quad Time: 20 min}
  \end{alertblock}
\end{frame}

% ------------------------------------------------------------
\subsection{Exercise 4D -- Line / Driver / Shift / Depot (5NF)}

\begin{frame}{Exercise 4D -- Line / Driver / Shift / Depot}
  \begin{block}{Starting Relation}
    $R($Line, Driver, Shift, Depot$)$
  \end{block}

  \begin{block}{Business Rules}
    \begin{itemize}
      \item Each driver can work on multiple lines.
      \item Each line operates in multiple shifts.
      \item Drivers are assigned to depots.
      \item The combination $(L, D, S, \text{Dp})$ is valid if and only if each pairwise assignment is valid \textbf{and} the combination is consistent.
    \end{itemize}
    This results in a \textbf{join dependency} $\bowtie\{(L,D),\,(L,S),\,(D,\text{Dp})\}$.
  \end{block}

  \begin{block}{Tasks}
    \begin{enumerate}
      \item Identify the join dependency.
      \item Decompose to \textbf{5NF} -- state each binary/ternary relation.
      \item Demonstrate that the original relation can be reconstructed by natural join.
    \end{enumerate}
  \end{block}

  \begin{alertblock}{Difficulty: Advanced \quad Time: 25 min}
  \end{alertblock}
\end{frame}

% ============================================================
\section{Canonical Cover \& Synthesis Algorithm}
% ============================================================

% ------------------------------------------------------------
\subsection{Exercise 5A}

\begin{frame}{Exercise 5A -- Canonical Cover \& Synthesis}
  \begin{block}{Given}
    $R = \{A, B, C, D\}$

    \vspace{4pt}
    $F = \{D \to AB,\quad CD \to B,\quad C \to B\}$
  \end{block}

  \begin{block}{Tasks}
    \begin{enumerate}
      \item Compute the \textbf{canonical cover} $F_c$:
        \begin{itemize}
          \item Decompose right-hand sides into single attributes.
          \item Eliminate extraneous attributes from left-hand sides.
          \item Eliminate redundant FDs.
        \end{itemize}
      \item Apply the \textbf{Synthesis Algorithm} to obtain a \textbf{3NF} decomposition.
      \item Check whether the result is:
        \begin{itemize}
          \item \textbf{Dependency-preserving}: can all FDs in $F$ be checked within a single table?
          \item \textbf{Lossless-join}: does the natural join of all tables reconstruct $R$?
        \end{itemize}
    \end{enumerate}
  \end{block}

  \begin{alertblock}{Difficulty: Intermediate \quad Time: 25 min}
  \end{alertblock}
\end{frame}

% ------------------------------------------------------------
\subsection{Exercise 5B}

\begin{frame}{Exercise 5B -- Canonical Cover \& Synthesis}
  \begin{block}{Given}
    $R = \{A, B, C, D, E\}$

    \vspace{4pt}
    $F = \{CD \to A,\quad D \to AE,\quad C \to BE,\quad ABD \to CE,\quad E \to A\}$
  \end{block}

  \begin{block}{Tasks}
    \begin{enumerate}
      \item Compute the \textbf{canonical cover} $F_c$.
      \item Determine all \textbf{candidate keys} of $R$.
      \item Apply the \textbf{Synthesis Algorithm} to obtain a \textbf{3NF} decomposition.
      \item Check whether the result is dependency-preserving and lossless-join.
    \end{enumerate}
  \end{block}

  \begin{block}{Reminder: Attribute Closure}
    To check if $X$ is a superkey: compute $X^+$ under $F$.\\
    If $X^+ = R$, then $X$ is a superkey.
  \end{block}

  \begin{alertblock}{Difficulty: Intermediate--Advanced \quad Time: 30 min}
  \end{alertblock}
\end{frame}

% ------------------------------------------------------------
\subsection{Exercise 5C}

\begin{frame}{Exercise 5C -- Canonical Cover \& Synthesis}
  \begin{block}{Given}
    $R = \{A, B, C, D, E, F\}$

    \vspace{4pt}
    $F = \{C \to DE,\quad ABC \to DF,\quad ABF \to CDE,\quad ABD \to CF,\quad CEF \to D\}$
  \end{block}

  \begin{block}{Tasks}
    \begin{enumerate}
      \item Compute the \textbf{canonical cover} $F_c$.
      \item Determine all \textbf{candidate keys} of $R$.
      \item Apply the \textbf{Synthesis Algorithm} to obtain a \textbf{3NF} decomposition.
      \item Check dependency-preservation and lossless-join.
      \item Optional: attempt \textbf{BCNF} decomposition and compare the result.
    \end{enumerate}
  \end{block}

  \begin{alertblock}{Difficulty: Advanced \quad Time: 35 min}
  \end{alertblock}
\end{frame}

% ============================================================
\section{Entity-Relationship Modelling (ERM)}
% ============================================================

% ------------------------------------------------------------
\subsection{Exercise 6A -- Airplane / Flight / Pilot (Simple)}

\begin{frame}{Exercise 6A -- Airplane / Flight / Pilot}
  \begin{block}{Description}
    \begin{itemize}
      \item An \textbf{Airplane} has: registration number, type, and capacity.
      \item A \textbf{Flight} has: flight number, departure airport, arrival airport, and date.
      \item A \textbf{Pilot} has: staff number, name, and license type.
      \item Each flight uses \textbf{exactly one} airplane.
      \item Each flight is operated by \textbf{exactly one} pilot.
      \item A pilot can operate \textbf{many} flights; an airplane can be used for \textbf{many} flights.
    \end{itemize}
  \end{block}

  \begin{block}{Task}
    Draw the \textbf{Entity-Relationship Diagram} in Chen's notation.
    \begin{itemize}
      \item Show all entities with their attributes (underline key attributes).
      \item Show all relationships with their names.
      \item Annotate all cardinalities (\texttt{1}, \texttt{N}, \texttt{M} notation or min--max).
    \end{itemize}
  \end{block}

  \begin{alertblock}{Difficulty: Simple \quad Time: 15 min}
  \end{alertblock}
\end{frame}

% ------------------------------------------------------------
\subsection{Exercise 6B -- Detailed Flight Schema (Intermediate)}

\begin{frame}{Exercise 6B -- Detailed Flight Schema}
  \begin{block}{Description}
    \begin{itemize}
      \item A \textbf{Flight} (identified by FlightNumber) consists of multiple \textbf{FlightLegs}.
      \item A \textbf{FlightLeg} connects one \textbf{Airport} to another (origin and destination).
      \item Each \textbf{Airport} has: IATA code, name, and city.
      \item A \textbf{DailyFlightLeg} is a scheduled instance of a FlightLeg on a specific date -- it is a \textbf{weak entity}.
      \item An \textbf{AircraftType} has: a name and seat configuration.
      \item Each DailyFlightLeg uses exactly one AircraftType.
    \end{itemize}
  \end{block}

  \begin{block}{Task}
    Draw the complete \textbf{Extended ERM} in Chen's notation.
    \begin{itemize}
      \item Identify and model \textbf{weak entities} with double rectangles and double diamonds.
      \item Annotate all cardinalities.
      \item Indicate all identifying relationships.
    \end{itemize}
  \end{block}

  \begin{alertblock}{Difficulty: Intermediate \quad Time: 25 min}
  \end{alertblock}
\end{frame}

% ------------------------------------------------------------
\subsection{Exercise 6C -- Crew with Specialization (Advanced)}

\begin{frame}{Exercise 6C -- Crew with Specialization}
  \begin{block}{Description}
    \begin{itemize}
      \item \textbf{Crew members} are either \textbf{Pilots} or \textbf{Flight Attendants} (total, disjoint specialization).
      \item All crew members share: staff number, name, \emph{languages spoken} (multivalued).
      \item \textbf{Pilots} additionally have: license number, license type, and a list of aircraft ratings (multivalued -- aircraft types they are certified for).
      \item \textbf{Flight Attendants} additionally have: seniority level.
      \item All crew members receive \textbf{emergency training}; each training record has a date and type.
    \end{itemize}
  \end{block}

  \begin{block}{Task}
    Draw the \textbf{Extended ERM} using generalization/specialization.
    \begin{itemize}
      \item Use the ISA triangle notation with disjoint/total markers.
      \item Show all attributes, including multivalued attributes (double ellipses).
      \item Show the \texttt{EmergencyTraining} relationship with its attributes.
      \item Annotate all cardinalities.
    \end{itemize}
  \end{block}

  \begin{alertblock}{Difficulty: Advanced \quad Time: 30 min}
  \end{alertblock}
\end{frame}

% ============================================================
\section{Extended ERM (EERM)}
% ============================================================

% ------------------------------------------------------------
\subsection{Exercise 7A -- Book / Chapter / Publisher / Author (Simple)}

\begin{frame}{Exercise 7A -- Book / Chapter / Publisher / Author}
  \begin{block}{Description}
    \begin{itemize}
      \item A \textbf{Book} has: title, ISBN, and publication year.
      \item Each Book belongs to exactly one \textbf{Publisher}; a Publisher has: name and address.
      \item Each Book has multiple \textbf{Chapters} -- a Chapter is a \textbf{weak entity} (exists only as part of a Book); it has a title and page numbers.
      \item \textbf{Authors} write Books (m:n); an Author can write multiple books; a Book can have multiple authors.
      \item An Author has: name and a list of previous publications (\textbf{multivalued attribute}).
    \end{itemize}
  \end{block}

  \begin{block}{Task}
    Draw the \textbf{Extended ERM}.
    \begin{itemize}
      \item Model the weak entity with double rectangle and double diamond.
      \item Model the multivalued attribute with double ellipse.
      \item Show the m:n relationship between Author and Book.
      \item Annotate all cardinalities.
    \end{itemize}
  \end{block}

  \begin{alertblock}{Difficulty: Simple \quad Time: 20 min}
  \end{alertblock}
\end{frame}

% ------------------------------------------------------------
\subsection{Exercise 7B -- Book with Overlapping Specialization (Intermediate)}

\begin{frame}{Exercise 7B -- Book with Overlapping Specialization}
  \begin{block}{Description}
    \begin{itemize}
      \item A \textbf{Book} is a weak entity with \textbf{Chapter} as an even weaker entity.
      \item \textbf{Contributors} are people who contribute to books; they can be \textbf{Authors}, \textbf{Editors}, or \textbf{both} -- this is an \textbf{overlapping, partial} specialization.
      \item \textbf{Authors} write Chapters (m:n); a Chapter can be written by multiple Authors.
      \item \textbf{Editors} edit entire Books (m:n).
      \item An Editor additionally has an editorial certification attribute.
    \end{itemize}
  \end{block}

  \begin{block}{Task}
    Draw the \textbf{Extended ERM}.
    \begin{itemize}
      \item Indicate \textbf{overlapping} (``O'') and \textbf{partial} (``P'') markers on the ISA triangle.
      \item Show the ``writes'' (Author--Chapter) and ``edits'' (Editor--Book) relationships.
      \item Annotate all cardinalities.
    \end{itemize}
  \end{block}

  \begin{alertblock}{Difficulty: Intermediate \quad Time: 25 min}
  \end{alertblock}
\end{frame}

% ------------------------------------------------------------
\subsection{Exercise 7C -- Complex Publication Schema (Advanced)}

\begin{frame}{Exercise 7C -- Complex Publication Schema}
  \begin{block}{Description}
    \begin{itemize}
      \item \textbf{Publication} is a supertype with subtypes \textbf{Book} and \textbf{Proceedings} (total, disjoint).
      \item A Book has \textbf{Editions} (weak entity); each Edition has \textbf{Chapters} (weak entity of Edition); each Chapter has \textbf{Sections} (weak entity of Chapter).
      \item A \textbf{Proceedings} belongs to a \textbf{Conference} (with name, year, location).
      \item A \textbf{Copy} is a physical instance of an Edition (weak entity of Edition).
      \item \textbf{Contributors} specialize into \textbf{Authors} and \textbf{Editors} (partial, overlapping).
    \end{itemize}
  \end{block}

  \begin{block}{Task}
    Draw the complete \textbf{Extended ERM}.
    \begin{itemize}
      \item Model all levels of nested weak entities.
      \item Show both specialization hierarchies.
      \item Annotate all cardinalities and identifying relationships.
    \end{itemize}
  \end{block}

  \begin{alertblock}{Difficulty: Advanced \quad Time: 35 min}
  \end{alertblock}
\end{frame}

% ============================================================
\section{ERM + Relational Schema Mapping}
% ============================================================

% ------------------------------------------------------------
\subsection{Exercise 8A -- Customers / Discounts / Items / Categories (Simple)}

\begin{frame}{Exercise 8A -- Customers / Discounts / Items / Categories}
  \begin{block}{Description}
    \begin{itemize}
      \item \textbf{Customers} can have \textbf{Discounts} (1:1 -- each customer has at most one active discount; a Discount has percentage and validity period).
      \item \textbf{Items} belong to \textbf{Categories} (m:n -- an item can be in multiple categories; a Category has name and description).
      \item An Item has: price, name, and description.
    \end{itemize}
  \end{block}

  \begin{block}{Tasks}
    \begin{enumerate}
      \item Draw the \textbf{EERM} with all entities, relationships, attributes, and cardinalities.
      \item \textbf{Map to a relational schema} applying all standard mapping rules:
        \begin{itemize}
          \item Strong entities $\to$ tables with PK.
          \item 1:1 relationship $\to$ FK in one of the participating tables.
          \item m:n relationship $\to$ separate junction table with composite PK.
        \end{itemize}
      \item Write \textbf{SQL \texttt{CREATE TABLE}} statements (DDL) for all tables, including all constraints.
    \end{enumerate}
  \end{block}

  \begin{alertblock}{Difficulty: Simple \quad Time: 25 min}
  \end{alertblock}
\end{frame}

% ------------------------------------------------------------
\subsection{Exercise 8B -- SAP Sales Document Schema (Intermediate)}

\begin{frame}{Exercise 8B -- SAP Sales Document Schema}
  \begin{block}{Description (simplified SAP SD module)}
    \begin{itemize}
      \item A \textbf{SalesDocument} has a \textbf{SalesDocumentType} (standard order, rush order, \ldots).
      \item Each SalesDocument goes through a \textbf{CreditCheckResult}.
      \item A SalesDocument produces a \textbf{BillingDocument} (\textbf{BillingType}: invoice, credit memo).
      \item Each SalesDocument generates a \textbf{DeliveryDocument} (\textbf{DeliveryType}: outbound delivery).
      \item Documents can carry a \textbf{Block} (delivery block, billing block, credit block).
    \end{itemize}
  \end{block}

  \begin{block}{Tasks}
    \begin{enumerate}
      \item Draw the \textbf{EERM} with all entities, type attributes, and cardinalities.
      \item Derive the complete \textbf{relational schema} with all FKs.
      \item Identify which relationships are 1:1, 1:N, or M:N.
    \end{enumerate}
  \end{block}

  \begin{alertblock}{Difficulty: Intermediate \quad Time: 30 min}
  \end{alertblock}
\end{frame}

% ------------------------------------------------------------
\subsection{Exercise 8C -- SAP Organizational Structure (Advanced)}

\begin{frame}{Exercise 8C -- SAP Organizational Structure}
  \begin{block}{Entities and Hierarchy}
    \small
    \begin{tabular}{@{}ll@{}}
      \textbf{Plant} & (PlantCode, Name, Address) \\
      \textbf{CompanyCode} & (CompanyCode, Name, Currency) \\
      \textbf{StorageLocation} & (StorageLocationCode, Name, PlantCode\textsuperscript{FK}) \\
      \textbf{PurchasingOrganization} & (PurchOrgCode, Name) \\
      \textbf{SalesOrganization} & (SalesOrgCode, Name, Currency, CompanyCode\textsuperscript{FK}) \\
      \textbf{DistributionChannel} & (DistChannelCode, Name) \\
      \textbf{Division} & (DivisionCode, Name) \\
    \end{tabular}
  \end{block}

  \begin{block}{Relationships}
    \begin{itemize}
      \item Plant \emph{belongs to} CompanyCode (N:1).
      \item StorageLocation \emph{belongs to} Plant (N:1).
      \item SalesOrganization \emph{belongs to} CompanyCode (N:1).
      \item SalesOrganization \emph{serves} Divisions \emph{through} DistributionChannels (ternary M:N:P).
    \end{itemize}
  \end{block}

  \begin{block}{Tasks}
    \begin{enumerate}
      \item Draw the complete \textbf{EERM}.
      \item Derive the complete \textbf{relational schema} with all FKs, including the ternary relationship.
    \end{enumerate}
  \end{block}

  \begin{alertblock}{Difficulty: Advanced \quad Time: 35 min}
  \end{alertblock}
\end{frame}

% ============================================================
\section{SQL -- SELECT}
% ============================================================

\begin{frame}{SQL Exercises -- Schema Overview}
  \begin{block}{Supermarket Schema}
    \small
    \begin{itemize}
      \item \textbf{Employee}(\underline{EmployeeID}, FirstName, LastName, Phone, DeptID\textsuperscript{FK}, HireDate)
      \item \textbf{Department}(\underline{DeptID}, Name, Floor)
      \item \textbf{Brand}(\underline{BrandID}, Name, Country, LegalForm)
      \item \textbf{Product}(\underline{EAN}, Name, Brand\textsuperscript{FK}, Price, AlcoholContent, DeptID\textsuperscript{FK})
      \item \textbf{Storage}(\underline{StorageID}, Floor, Capacity)
      \item \textbf{StorageProduct}(\underline{StorageID\textsuperscript{FK}, EAN\textsuperscript{FK}}, Quantity)
    \end{itemize}
  \end{block}

  \begin{block}{Music Schema}
    \small
    \begin{itemize}
      \item \textbf{Artist}(\underline{ArtistID}, Name, Country, DebutYear)
      \item \textbf{Album}(\underline{AlbumID}, Title, ArtistID\textsuperscript{FK}, Genre, ReleaseYear, RecordLabel)
      \item \textbf{Track}(\underline{TrackID}, Title, AlbumID\textsuperscript{FK}, Duration, FeaturedArtistID\textsuperscript{FK} \textit{nullable})
    \end{itemize}
  \end{block}

  These two schemas are used throughout all SQL exercises (Sections 9--12).
\end{frame}

% ------------------------------------------------------------
\subsection{Exercise 9A -- Supermarket -- Basic SELECT}

\begin{frame}[fragile]{Exercise 9A -- Supermarket -- Basic SELECT}
  \begin{block}{Tasks}
    \begin{enumerate}
      \item Show all employee names and phone numbers.
      \item Show all brand names and their legal forms.
      \item Show the first 5 products ordered by name (alphabetically).
      \item List all products with alcohol content $> 0$.
      \item Show distinct EAN codes aliased as \texttt{barcode}.
      \item Show the first 5 employees (use \texttt{LIMIT}).
    \end{enumerate}
  \end{block}

  \begin{block}{Example}
\begin{lstlisting}[style=sqlstyle]
-- Task 1 example
SELECT FirstName, LastName, Phone
FROM Employee;
\end{lstlisting}
  \end{block}

  \begin{alertblock}{Difficulty: Simple \quad Time: 10 min}
  \end{alertblock}
\end{frame}

% ------------------------------------------------------------
\subsection{Exercise 9B -- Music -- Basic SELECT}

\begin{frame}[fragile]{Exercise 9B -- Music -- Basic SELECT}
  \begin{block}{Tasks}
    \begin{enumerate}
      \item Show all album names and genres.
      \item Show the first 3 record labels (distinct, alphabetically).
      \item Show track titles aliased as \texttt{song}.
      \item Show distinct countries of artists.
      \item Show the first 5 artists.
    \end{enumerate}
  \end{block}

  \begin{block}{Hints}
\begin{lstlisting}[style=sqlstyle]
-- Alias a column
SELECT Title AS song FROM Track;

-- Distinct values
SELECT DISTINCT Country FROM Artist;

-- Limit results
SELECT Name FROM Artist LIMIT 5;
\end{lstlisting}
  \end{block}

  \begin{alertblock}{Difficulty: Simple \quad Time: 10 min}
  \end{alertblock}
\end{frame}

% ------------------------------------------------------------
\subsection{Exercise 9C -- Supermarket -- WHERE Clauses}

\begin{frame}[fragile]{Exercise 9C -- Supermarket -- WHERE Clauses}
  \begin{block}{Tasks}
    \begin{enumerate}
      \item List all brands from Germany.
      \item List products not stored anywhere (no entry in StorageProduct).
      \item List employees whose last name starts with \texttt{'N'}.
      \item List products with alcohol content between 1\% and 20\%.
      \item List brands that are a \texttt{GmbH} in Germany but \textbf{not} in Hamburg.
    \end{enumerate}
  \end{block}

  \begin{block}{Useful Operators}
\begin{lstlisting}[style=sqlstyle]
-- Pattern matching
WHERE LastName LIKE 'N%'

-- Range filter
WHERE AlcoholContent BETWEEN 1 AND 20

-- Negation + multiple conditions
WHERE Country = 'Germany'
  AND LegalForm = 'GmbH'
  AND City <> 'Hamburg'
\end{lstlisting}
  \end{block}

  \begin{alertblock}{Difficulty: Simple \quad Time: 15 min}
  \end{alertblock}
\end{frame}

% ------------------------------------------------------------
\subsection{Exercise 9D -- Music -- WHERE Clauses}

\begin{frame}[fragile]{Exercise 9D -- Music -- WHERE Clauses}
  \begin{block}{Tasks}
    \begin{enumerate}
      \item List all distinct genres in the Album table.
      \item List songs that have a featured artist (\texttt{FeaturedArtistID} is not NULL).
      \item List songs with a title starting with \texttt{'L'}.
      \item List songs with duration between 2 and 3 minutes (120 and 180 seconds).
      \item List all artists from the USA.
    \end{enumerate}
  \end{block}

  \begin{block}{Hint}
\begin{lstlisting}[style=sqlstyle]
-- Check for NULL
WHERE FeaturedArtistID IS NOT NULL

-- Duration in seconds
WHERE Duration BETWEEN 120 AND 180
\end{lstlisting}
  \end{block}

  \begin{alertblock}{Difficulty: Simple \quad Time: 15 min}
  \end{alertblock}
\end{frame}

% ------------------------------------------------------------
\subsection{Exercise 9E -- Supermarket -- Aggregation}

\begin{frame}[fragile]{Exercise 9E -- Supermarket -- Aggregation}
  \begin{block}{Tasks}
    \begin{enumerate}
      \item Count the number of \textbf{distinct} countries of origin for brands.
      \item Create price categories using \texttt{CASE WHEN}:
        \begin{itemize}
          \item Price $< 1.00$ \euro{} $\Rightarrow$ \texttt{'cheap'}
          \item $1.00$--$5.00$ \euro{} $\Rightarrow$ \texttt{'medium'}
          \item Price $> 5.00$ \euro{} $\Rightarrow$ \texttt{'expensive'}
        \end{itemize}
      \item Show the number of products per country and city, ordered by count descending.
      \item Show the \textbf{average price} per department, excluding departments 1 and 2.
      \item Show departments with at least \textbf{5 distinct barcodes} using \texttt{HAVING}.
    \end{enumerate}
  \end{block}

  \begin{block}{Hint for Task 2}
\begin{lstlisting}[style=sqlstyle]
SELECT Name,
  CASE WHEN Price < 1.00 THEN 'cheap'
       WHEN Price <= 5.00 THEN 'medium'
       ELSE 'expensive'
  END AS PriceCategory
FROM Product;
\end{lstlisting}
  \end{block}

  \begin{alertblock}{Difficulty: Intermediate \quad Time: 20 min}
  \end{alertblock}
\end{frame}

% ------------------------------------------------------------
\subsection{Exercise 9F -- Music -- Aggregation}

\begin{frame}[fragile]{Exercise 9F -- Music -- Aggregation}
  \begin{block}{Tasks}
    \begin{enumerate}
      \item Show all distinct countries of record labels (i.e.\ all countries artists come from).
      \item Categorize songs by duration using \texttt{CASE WHEN}:
        \begin{itemize}
          \item $< 120$ s $\Rightarrow$ \texttt{'short'}, \quad $120$--$240$ s $\Rightarrow$ \texttt{'medium'}, \quad $> 240$ s $\Rightarrow$ \texttt{'long'}
        \end{itemize}
      \item Count artists per country, ordered by count descending.
      \item Show the total duration per album in seconds, excluding the \texttt{'Christmas'} genre.
      \item Show artists who have at least \textbf{5 songs} using \texttt{HAVING}.
    \end{enumerate}
  \end{block}

  \begin{block}{Hint for Task 5}
\begin{lstlisting}[style=sqlstyle]
SELECT ArtistID, COUNT(*) AS SongCount
FROM Track
  JOIN Album ON Track.AlbumID = Album.AlbumID
GROUP BY ArtistID
HAVING COUNT(*) >= 5;
\end{lstlisting}
  \end{block}

  \begin{alertblock}{Difficulty: Intermediate \quad Time: 20 min}
  \end{alertblock}
\end{frame}

% ============================================================
\section{SQL -- JOINs}
% ============================================================

% ------------------------------------------------------------
\subsection{Exercise 10A -- Music -- JOINs (Simple)}

\begin{frame}[fragile]{Exercise 10A -- Music -- JOINs (Simple)}
  \begin{block}{Tasks}
    \begin{enumerate}
      \item List all songs in the \texttt{'Pop'} genre (JOIN Track with Album).
      \item List artists who have songs on Taylor Swift albums.
      \item List all albums from an Australian artist.
      \item Show each album with its artist name.
      \item Show track names with album title \textbf{and} artist name (three-table join).
    \end{enumerate}
  \end{block}

  \begin{block}{Reminder: JOIN syntax}
\begin{lstlisting}[style=sqlstyle]
-- INNER JOIN
SELECT t.Title, a.Title AS Album
FROM Track t
  JOIN Album a ON t.AlbumID = a.AlbumID;

-- THREE-TABLE JOIN
SELECT t.Title, al.Title, ar.Name
FROM Track t
  JOIN Album al ON t.AlbumID = al.AlbumID
  JOIN Artist ar ON al.ArtistID = ar.ArtistID;
\end{lstlisting}
  \end{block}

  \begin{alertblock}{Difficulty: Simple \quad Time: 15 min}
  \end{alertblock}
\end{frame}

% ------------------------------------------------------------
\subsection{Exercise 10B -- Music -- JOINs (Intermediate)}

\begin{frame}[fragile]{Exercise 10B -- Music -- JOINs (Intermediate)}
  \begin{block}{Tasks}
    \begin{enumerate}
      \item List all songs \textbf{not} in the \texttt{'Christmas'} genre.
      \item List artists who have songs \textbf{shorter than 3 minutes} (180 s).
      \item List artists who have \textbf{more than 5 albums}.
      \item Show artists with their \textbf{total number of tracks}.
      \item List albums that have \textbf{no tracks} (use \texttt{LEFT JOIN} + \texttt{IS NULL}).
    \end{enumerate}
  \end{block}

  \begin{block}{Hint for Task 5 -- Finding albums with no tracks}
\begin{lstlisting}[style=sqlstyle]
SELECT al.Title
FROM Album al
  LEFT JOIN Track t ON al.AlbumID = t.AlbumID
WHERE t.TrackID IS NULL;
\end{lstlisting}
  \end{block}

  \begin{alertblock}{Difficulty: Intermediate \quad Time: 20 min}
  \end{alertblock}
\end{frame}

% ------------------------------------------------------------
\subsection{Exercise 10C -- Supermarket -- JOINs (Intermediate)}

\begin{frame}[fragile]{Exercise 10C -- Supermarket -- JOINs}
  \begin{block}{Tasks}
    \begin{enumerate}
      \item List all products from \textbf{France} (join Brand and Product).
      \item List products that are in storage but have \textbf{no department} assigned.
      \item List \textbf{empty storage rooms} (no products stored there).
      \item List \textbf{brands with no products}.
      \item List employees who work in departments that have at least \textbf{3 products from the USA}.
    \end{enumerate}
  \end{block}

  \begin{block}{Hint for Task 3}
\begin{lstlisting}[style=sqlstyle]
-- Empty storage rooms: LEFT JOIN + IS NULL
SELECT s.StorageID, s.Floor
FROM Storage s
  LEFT JOIN StorageProduct sp
         ON s.StorageID = sp.StorageID
WHERE sp.EAN IS NULL;
\end{lstlisting}
  \end{block}

  \begin{alertblock}{Difficulty: Intermediate \quad Time: 25 min}
  \end{alertblock}
\end{frame}

% ============================================================
\section{SQL -- INSERT / UPDATE / DELETE}
% ============================================================

% ------------------------------------------------------------
\subsection{Exercise 11A -- INSERT -- Supermarket}

\begin{frame}[fragile]{Exercise 11A -- INSERT -- Supermarket}
  \begin{block}{Tasks}
    \begin{enumerate}
      \item Insert a new department \texttt{'Bakery'} on floor 2 into the Department table.
      \item Insert a new employee \texttt{'John Smith'} who works in the Bakery department (use the ID you just inserted).
    \end{enumerate}
  \end{block}

  \begin{block}{Syntax Reminder}
\begin{lstlisting}[style=sqlstyle]
-- Full insert
INSERT INTO table_name (col1, col2, col3)
VALUES (val1, val2, val3);

-- Insert using LAST_INSERT_ID() for FK
INSERT INTO Employee (FirstName, LastName, DeptID)
VALUES ('John', 'Smith', LAST_INSERT_ID());
\end{lstlisting}
  \end{block}

  \begin{alertblock}{Difficulty: Simple \quad Time: 10 min}
  \end{alertblock}
\end{frame}

% ------------------------------------------------------------
\subsection{Exercise 11B -- UPDATE -- Supermarket}

\begin{frame}[fragile]{Exercise 11B -- UPDATE -- Supermarket}
  \begin{block}{Tasks}
    \begin{enumerate}
      \item Rename the department \texttt{'Non-food'} to \texttt{'Household'}.
      \item Update the price of the product \texttt{'Coca Cola Zero'} to \texttt{1.29}.
    \end{enumerate}
  \end{block}

  \begin{block}{Syntax Reminder}
\begin{lstlisting}[style=sqlstyle]
UPDATE table_name
SET column = new_value
WHERE condition;
\end{lstlisting}
  \end{block}

  \begin{alertblock}{Always use WHERE with UPDATE!}
    Without a \texttt{WHERE} clause, \textbf{every row} in the table will be updated.
    Check your condition carefully before executing.
  \end{alertblock}

  \begin{alertblock}{Difficulty: Simple \quad Time: 10 min}
  \end{alertblock}
\end{frame}

% ------------------------------------------------------------
\subsection{Exercise 11C -- DELETE -- Supermarket}

\begin{frame}[fragile]{Exercise 11C -- DELETE -- Supermarket}
  \begin{block}{Tasks}
    \begin{enumerate}
      \item Delete the department \texttt{'Flowers and Gardening'}.
      \item Delete the employee with last name \texttt{'Major'}.
    \end{enumerate}
  \end{block}

  \begin{block}{Syntax Reminder}
\begin{lstlisting}[style=sqlstyle]
DELETE FROM table_name
WHERE condition;
\end{lstlisting}
  \end{block}

  \begin{alertblock}{Referential Integrity Warning}
    Deleting a department that is referenced by employees will \textbf{fail} unless you:
    \begin{itemize}
      \item Delete the employees first (manual cascade).
      \item Have defined \texttt{ON DELETE CASCADE} on the FK.
      \item Set employee \texttt{DeptID} to \texttt{NULL} first (\texttt{ON DELETE SET NULL}).
    \end{itemize}
  \end{alertblock}

  \begin{alertblock}{Difficulty: Simple \quad Time: 10 min}
  \end{alertblock}
\end{frame}

% ------------------------------------------------------------
\subsection{Exercise 11D -- INSERT -- Music}

\begin{frame}[fragile]{Exercise 11D -- INSERT -- Music}
  \begin{block}{Tasks}
    \begin{enumerate}
      \item Insert a new \textbf{Artist} of your choice into the Artist table.
      \item Insert one of their \textbf{Albums} into the Album table (referencing the artist you just created).
      \item Insert at least \textbf{3 Tracks} from that album into the Track table.
    \end{enumerate}
  \end{block}

  \begin{block}{Remember to respect FK constraints}
\begin{lstlisting}[style=sqlstyle]
-- Step 1: insert artist
INSERT INTO Artist (ArtistID, Name, Country, DebutYear)
VALUES (999, 'My Artist', 'Germany', 2020);

-- Step 2: insert album referencing that artist
INSERT INTO Album (AlbumID, Title, ArtistID, Genre,
                   ReleaseYear, RecordLabel)
VALUES (888, 'My Album', 999, 'Pop', 2022, 'My Label');
\end{lstlisting}
  \end{block}

  \begin{alertblock}{Difficulty: Simple \quad Time: 15 min}
  \end{alertblock}
\end{frame}

% ------------------------------------------------------------
\subsection{Exercise 11E -- UPDATE -- Music}

\begin{frame}[fragile]{Exercise 11E -- UPDATE -- Music}
  \begin{block}{Tasks}
    \begin{enumerate}
      \item Change the genre of album \texttt{'Fearless'} to \texttt{'Country Pop'}.
      \item Change the record label of album \texttt{'Red'} to \texttt{'Republic Records'}.
      \item Revert the change in task 2: set the record label of \texttt{'Red'} back to its original value.
    \end{enumerate}
  \end{block}

  \begin{block}{Useful pattern -- check before updating}
\begin{lstlisting}[style=sqlstyle]
-- Always SELECT first to verify your WHERE clause
SELECT AlbumID, Title, RecordLabel
FROM Album
WHERE Title = 'Red';

-- Then UPDATE
UPDATE Album
SET RecordLabel = 'Republic Records'
WHERE Title = 'Red';
\end{lstlisting}
  \end{block}

  \begin{alertblock}{Difficulty: Simple \quad Time: 10 min}
  \end{alertblock}
\end{frame}

% ------------------------------------------------------------
\subsection{Exercise 11F -- DELETE -- Music}

\begin{frame}[fragile]{Exercise 11F -- DELETE -- Music}
  \begin{block}{Tasks}
    \begin{enumerate}
      \item Delete all tracks that \textbf{feature} the artist \texttt{'Blackpink'} (i.e.\ where \texttt{FeaturedArtistID} refers to Blackpink).
      \item Delete the artist \texttt{'Ice Spice'} from the Artist table.\\
            \textit{Note: handle FK constraints -- you may need to nullify or delete referencing rows first.}
    \end{enumerate}
  \end{block}

  \begin{block}{Handling FK constraints before DELETE}
\begin{lstlisting}[style=sqlstyle]
-- Option A: nullify the featured artist reference first
UPDATE Track
SET FeaturedArtistID = NULL
WHERE FeaturedArtistID = (
  SELECT ArtistID FROM Artist WHERE Name = 'Ice Spice'
);

-- Then delete the artist
DELETE FROM Artist WHERE Name = 'Ice Spice';
\end{lstlisting}
  \end{block}

  \begin{alertblock}{Difficulty: Simple--Intermediate \quad Time: 15 min}
  \end{alertblock}
\end{frame}

% ============================================================
\section{SQL -- DDL}
% ============================================================

% ------------------------------------------------------------
\subsection{Exercise 12A -- CREATE / ALTER -- Music}

\begin{frame}[fragile]{Exercise 12A -- CREATE / ALTER -- Music}
  \begin{block}{Tasks}
    \begin{enumerate}
      \item Create a new table \textbf{Concert}(\underline{ConcertID}, ArtistID, Venue, City, Date, Capacity).
      \item Add a \textbf{foreign key} constraint from \texttt{Concert.ArtistID} to \texttt{Artist.ArtistID}.
      \item Add a column \texttt{Website} (\texttt{VARCHAR(200)}, nullable) to the \texttt{Artist} table.
      \item Change the \texttt{Genre} column in \texttt{Album} to \texttt{VARCHAR(100)}.
      \item Drop the \texttt{Country} column from \texttt{Artist}.
    \end{enumerate}
  \end{block}

  \begin{block}{Syntax Reminder}
\begin{lstlisting}[style=sqlstyle]
CREATE TABLE Concert (
  ConcertID INT PRIMARY KEY AUTO_INCREMENT,
  ArtistID  INT NOT NULL,
  Venue     VARCHAR(200),
  City      VARCHAR(100),
  Date      DATE,
  Capacity  INT
);

ALTER TABLE Concert
  ADD CONSTRAINT fk_concert_artist
    FOREIGN KEY (ArtistID) REFERENCES Artist(ArtistID);

ALTER TABLE Artist ADD COLUMN Website VARCHAR(200);
ALTER TABLE Album  MODIFY COLUMN Genre VARCHAR(100);
ALTER TABLE Artist DROP COLUMN Country;
\end{lstlisting}
  \end{block}

  \begin{alertblock}{Difficulty: Simple \quad Time: 15 min}
  \end{alertblock}
\end{frame}

% ------------------------------------------------------------
\subsection{Exercise 12B -- CREATE / ALTER -- Supermarket}

\begin{frame}[fragile]{Exercise 12B -- CREATE / ALTER -- Supermarket}
  \begin{block}{Tasks}
    \begin{enumerate}
      \item Add columns \texttt{Salary} (\texttt{DECIMAL(10,2)}) and \texttt{HireDate} (\texttt{DATE}) to the \texttt{Employee} table.
      \item Create a \textbf{Manager} table: (\underline{ManagerID}, EmployeeID\textsuperscript{FK}, Salary, StartDate).
      \item Add a column \texttt{ManagerID} (nullable FK) to the \texttt{Department} table -- a department can have one manager.
    \end{enumerate}
  \end{block}

  \begin{block}{Syntax Reminder}
\begin{lstlisting}[style=sqlstyle]
ALTER TABLE Employee
  ADD COLUMN Salary   DECIMAL(10,2),
  ADD COLUMN HireDate DATE;

CREATE TABLE Manager (
  ManagerID  INT PRIMARY KEY AUTO_INCREMENT,
  EmployeeID INT NOT NULL,
  Salary     DECIMAL(10,2),
  StartDate  DATE,
  CONSTRAINT fk_mgr_emp
    FOREIGN KEY (EmployeeID) REFERENCES Employee(EmployeeID)
);

ALTER TABLE Department
  ADD COLUMN ManagerID INT NULL,
  ADD CONSTRAINT fk_dept_mgr
    FOREIGN KEY (ManagerID) REFERENCES Manager(ManagerID);
\end{lstlisting}
  \end{block}

  \begin{alertblock}{Difficulty: Simple \quad Time: 15 min}
  \end{alertblock}
\end{frame}

% ------------------------------------------------------------
\subsection{Exercise 12C -- CREATE DATABASE from Scratch}

\begin{frame}[fragile]{Exercise 12C -- Cinema Management Database}
  \begin{block}{Tasks}
    \begin{enumerate}
      \item Create a database named \texttt{cinema\_management}.
      \item Create table \textbf{Movie}(\underline{MovieID} AUTO\_INCREMENT, Title, Genre, Duration, Rating DECIMAL, ReleaseYear).
      \item Create table \textbf{Cinema}(\underline{CinemaID}, Name, City, Screens INT).
      \item Create table \textbf{Screening}(\underline{ScreeningID}, MovieID\textsuperscript{FK}, CinemaID\textsuperscript{FK}, ScreeningDate DATE, StartTime TIME, Seats INT, Price DECIMAL).
      \item Add all appropriate constraints (\texttt{NOT NULL}, \texttt{CHECK}, \texttt{DEFAULT}).
    \end{enumerate}
  \end{block}

  \begin{block}{Partial Skeleton}
\begin{lstlisting}[style=sqlstyle]
CREATE DATABASE cinema_management;
USE cinema_management;

CREATE TABLE Movie (
  MovieID     INT PRIMARY KEY AUTO_INCREMENT,
  Title       VARCHAR(255) NOT NULL,
  Genre       VARCHAR(100),
  Duration    INT CHECK (Duration > 0),
  Rating      DECIMAL(3,1),
  ReleaseYear INT
);
-- continue with Cinema and Screening ...
\end{lstlisting}
  \end{block}

  \begin{alertblock}{Difficulty: Intermediate \quad Time: 25 min}
  \end{alertblock}
\end{frame}

% ============================================================
\section{Group Exercise}
% ============================================================

% ------------------------------------------------------------
\subsection{Exercise 13 -- One-Pager Presentation}

\begin{frame}{Exercise 13 -- One-Pager Presentation}
  \begin{block}{Instructions}
    Form groups of \textbf{2--3 students}. Each group picks \textbf{one} of the topics below and prepares a \textbf{5-minute presentation}.
  \end{block}

  \begin{block}{Option A: Concurrency Problems}
    Explain \textbf{Lost Update}, \textbf{Dirty Read}, and \textbf{Phantom Read} with concrete, realistic examples. What transaction isolation level prevents each problem?
  \end{block}

  \begin{block}{Option B: CAP Theorem}
    Explain the three properties (\textbf{Consistency}, \textbf{Availability}, \textbf{Partition Tolerance}) and the trade-offs. Give real-world database examples for each combination (CA, CP, AP).
  \end{block}

  \begin{block}{Option C: BASE vs.\ ACID}
    Compare the two consistency models. Explain when each is appropriate. Give concrete examples of systems that use ACID and systems that use BASE.
  \end{block}
\end{frame}

\begin{frame}{Exercise 13 -- Requirements}
  \begin{block}{Deliverables}
    \begin{enumerate}
      \item \textbf{One-pager}: maximum one A4 page to serve as handout / visual aid for classmates.
        \begin{itemize}
          \item Must be self-contained: someone who hasn't studied the topic should understand the key ideas.
          \item Include at least one concrete example.
        \end{itemize}
      \item \textbf{Presentation}: exactly 5 minutes (use a timer).
        \begin{itemize}
          \item All group members must speak.
          \item Explain the concept, not just list definitions.
        \end{itemize}
    \end{enumerate}
  \end{block}

  \begin{block}{Evaluation Criteria}
    \begin{itemize}
      \item Correctness and completeness of the content.
      \item Clarity of explanation to peers.
      \item Quality and usefulness of the one-pager.
      \item Adherence to the 5-minute time limit.
    \end{itemize}
  \end{block}

  \begin{alertblock}{Preparation time: 20 minutes \quad Presentations: in class}
  \end{alertblock}
\end{frame}

\end{document}
